% =====================================================================
\section{Interactive Graphs: Desmos and Python}\label{app:interactive}
% =====================================================================

\subsection{Desmos}

The slide's graph \extlink{https://www.desmos.com/calculator/cho2h9ztfa}{desmos.com/calculator/cho2h9ztfa}
draws the gradient field of Example~\ref{ex:7.2} as arrows on a grid. To draw
your own planar field $\vF=\ip{M,N}$ in the
\extlink{\desmosTwoD}{Desmos calculator}:
\begin{enumerate}
\item Define the components, e.g. \texttt{M(x,y)=-y} and \texttt{N(x,y)=x}.
\item Make a grid of base points: \texttt{X=[-3,-2,...,3]}, \texttt{Y=[-3,-2,...,3]}.
\item Draw the arrows with \texttt{vector((X,Y),(X+0.3M(X,Y),Y+0.3N(X,Y)))}
      (repeat with \texttt{Y} fixed or use a list of points); the factor $0.3$
      scales the arrows.
\end{enumerate}
Adding the level curves \texttt{x\^{}2y-y\^{}3=k} with a slider $k$ shows that a
gradient field is perpendicular to the level curves of its potential.

\subsection{Python}

All figures are produced by \texttt{figures/make\_figures.py}. The core of a
vector-field plot is two lines of \texttt{matplotlib}:
\begin{center}
\begin{minipage}{0.85\linewidth}
\small
\texttt{X, Y = np.meshgrid(np.linspace(-2, 2, 11), np.linspace(-2, 2, 11))}\\
\texttt{plt.quiver(X, Y, -Y, X)\quad\# the field <-y, x>}
\end{minipage}
\end{center}
Replace \texttt{-Y, X} by any components $M(X,Y)$, $N(X,Y)$; for fields in space use
\texttt{ax.quiver(X, Y, Z, M, N, P)} on a 3D axis. In Jupyter or Google Colab,
\texttt{\%matplotlib widget} makes the 3D plots rotatable.
